% authority: science-draft
% object_id: TEX-V22-P4-T02-B2-POPULATED-INSTANCE-SMUGGLING-AUDIT-V1
% task_id: RT-20260809-026
\documentclass[11pt]{article}
\usepackage[margin=1in]{geometry}
\usepackage{amsmath,amssymb,amsthm,booktabs,array,longtable,enumitem,microtype}
\usepackage[hidelinks]{hyperref}
\setlength{\emergencystretch}{3em}
\newtheorem{definition}{Definition}
\newtheorem{proposition}{Proposition}
\newtheorem{lemma}{Lemma}
\newtheorem{corollary}{Corollary}
\newcommand{\R}{\mathbb{R}}
\newcommand{\cI}{\mathcal{I}}
\newcommand{\cC}{\mathcal{C}}
\newcommand{\cS}{\mathcal{S}}
\newcommand{\rad}{\operatorname{rad}}
\newcommand{\Char}{\operatorname{Char}}
\newcommand{\id}{\operatorname{id}}
\title{V22 P4--T02 B2 Populated-Instance Smuggling Audit\\
\large Repair-Required Verdict for the Equipped-Chain Descriptor Attempt}
\author{Smuggling Auditor task overlay RT-20260809-026}
\date{2026-08-10}

\begin{document}
\maketitle

\begin{abstract}
This draft/control audit examines the exact populated descriptor candidate
\begin{center}
\small\texttt{CAND-V22-B2-P7-COMMON-PRINCIPAL-LIFT-V1}.
\end{center}
It distinguishes algebraic
correctness inside the proposed equipment from source derivation of that
equipment.  The constant lift, sampling maps, identity quotient groupoids,
unit-one cocycle, and common-ideal calculation are internally correct.
Nevertheless, the populated attempt does not earn descriptor-instance credit.
Five defects remain: incomplete sector coverage; a presentation-dependent
neighborhood norm; insertion of one shared transport field into every sector
operator; a variation class that excludes arbitrarily small sector-split
perturbations; and an untyped operational composition.  The exact disposition
is \nolinkurl{repair_required_no_instance_credit}.  This is neither a global
no-go theorem nor a rejection of the candidate family.  No adequacy
reevaluation, B2 activation, P4--T03 unlock, target metric, or protected
authority is supplied.
\end{abstract}

\section{Authority and scope}

The audited subject is the RT-20260809-023 proposal-only witness, including
its population matrix, source-factorization provenance, explicit-unit
cocycle, and D7 separation obstruction.  The governing descriptor schema is
the RT-20260809-021 B2 continuum-lift descriptor specification, as audited at
schema level by RT-20260809-022.  The physical input is the protected P7
conditional-input package and its P7--T03, P7--T04, and P7--T05 records.

The present task asks only whether the \emph{populated} source fields and maps
obtain their advertised status without importing or preloading the desired
commonity.  The audit may:
\begin{itemize}[leftmargin=2em]
  \item verify source-fiber typing, algebra, covariance, and operational types;
  \item construct local countermodels to a claimed robustness or invariance;
  \item separate direct target import from source-side goal-property preload;
  \item return a bounded repair contract or a scoped obstruction.
\end{itemize}
It may not evaluate D7 adequacy, activate B2, execute P4--T03, adopt an
ontology law, or promote a physical causal geometry.

\paragraph{Protected P7 status.}
The P7 package is an adopted physical-matter conditional input, not a
first-principles derivation.  Its finite records equip three protocol controls
(Rod, Signal, Detector) on a common address chain.  Those records do not state
that the three controls exhaust every admitted P7 matter sector, nor do they
derive a constructor equipping every such sector.  P7 guarded labels are
syntax; they are not by themselves physical devices, PDE fields, or causal
cones.

\section{Audited populated witness}

On the proposal-only patch
\[
 U=(-1,1)^4\subset \R^4,
 \qquad \tau=\frac{\partial}{\partial \xi^0},
\]
the attempt uses three sector labels
\[
 \cS_0=\{R,S,D\},
\]
constant-copy lifts and finite-grid sampling, identity quotient groupoids, and
affine equations
\begin{equation}
 E_s(u,r)=\tau(u)+(I-M_B)(u-Lr),\qquad s\in\cS_0.
 \label{eq:operator}
\end{equation}
Here the zero-order carrier matrix is
\[
 M_B=
 \begin{pmatrix}
 0&0&1\\
 1&0&0\\
 0&1&0
 \end{pmatrix},
\]
the cyclic permutation matrix for the equipped three-address chain.  The
principal symbol of every displayed operator is
\begin{equation}
 \sigma_s(k)=k(\tau)I_3,\qquad
 \rad\cI_s=\langle k(\tau)\rangle.
 \label{eq:ideal}
\end{equation}
The compatibility set is then declared by equality with the predeclared ideal
\[
 \cC_\tau=\{s:\rad\cI_s=\langle k(\tau)\rangle\}.
\]
The allowed variations use one common perturbed field \(\tau'\) in all three
operators.  The operational paragraph proposes finite samples
\[
 S_{s,0}:\Gamma(U,\R^3)\longrightarrow(\R^3)^{G_0}
\]
followed by the P7--T03 response rule.

\section{Audit test: weak factorization versus anchored derivation}

\begin{definition}[Weak source-fiber factorization]
A populated descriptor datum \(d\) weakly factors through source data when
there exist a source object \(x\) and a map \(F\) such that \(d=F(x)\), and
the codomain of \(F\) contains no target metric, target atlas, or benchmark
output.
\end{definition}

\begin{definition}[Anchored source derivation]
A datum is anchored when, in addition to weak factorization,
\begin{enumerate}[label=(A\arabic*),leftmargin=2.5em]
  \item the source theory selects the domain objects and the map's type;
  \item every declared physical sector is covered or the scope is explicitly
        restricted to a named subfamily;
  \item auxiliary topology, norm, and variation class are intrinsic or their
        presentation group is declared and proved harmless;
  \item a desired cross-sector property is concluded from sector laws, rather
        than inserted as a shared primitive; and
  \item each operational composition is type-correct.
\end{enumerate}
\end{definition}

\begin{proposition}[Weak factorization is insufficient]
Every populated component of the RT-023 attempt can be written as a function
of the proposed patch, grids, transport field, carrier matrix, lift, and
sampling data.  This proves weak source-fiber factorization only.  It does not
prove anchored source derivation because the finite P7 records do not select
the patch, universal transport field, norm, restricted variation class, or
operational bridge.
\end{proposition}

\begin{proof}
The component formulas refer only to the named proposal-only source objects,
so weak factorization holds by inspection.  Conditions (A2)--(A5), however,
are independent typing and invariance burdens.  Sections 4--8 exhibit one
failed witness for each.  None is repaired by the absence of a target metric
symbol in \eqref{eq:operator}.
\end{proof}

This distinction is the core audit result: direct target import is absent, but
the desired common-leading property is preloaded by a new source primitive.

\section{Finding F1: three controls do not cover every P7 sector}

Descriptor component D0 requires a declared family of physical sectors.  The
attempt substitutes the three equipped protocol controls
\(R,S,D\) for the full admitted P7 matter-sector family.  P7--T04 records that
the wider protocol suite does not carry a common bijective relation, while
P7--T05 expressly limits the common-address witness to three controls and says
that not every admitted sector is equipped.

\begin{proposition}[Coverage obstruction]
The inclusion
\[
 \{R,S,D\}\hookrightarrow \cS_{P7}
\]
does not prove equality.  Therefore a descriptor quantified over every
declared matter sector cannot obtain D0 credit from the three-chain witness.
\end{proposition}

\begin{proof}
The cited P7 record supplies only a restricted common-address subfamily and
contains no exhaustivity theorem or all-sector equipment constructor.
Interpreting controls as an exhaustive sector taxonomy adds a new premise.
\end{proof}

This is a repairable typing defect: either supply an exact all-sector
constructor, or restate the candidate honestly as an equipped-subfamily
control whose result cannot yet discharge the full B2 burden.

\section{Finding F2: the neighborhood norm is presentation-dependent}

The populated attempt uses a componentwise \(C^1\) sup topology and a
threshold such as \(\lVert \delta\tau\rVert_\infty<1/2\) in the
\(\xi\)-presentation.

\begin{lemma}[Coordinate-rescaling witness]
Let \(\eta^0=2\xi^0\) and \(\eta^i=\xi^i\) for \(i=1,2,3\).  If
\[
 \delta\tau=a\frac{\partial}{\partial \xi^0},
\]
then its time component in the \(\eta\)-basis is \(2a\).  Consequently the
componentwise threshold \(\lvert a\rvert<1/2\) is not invariant under this
source-patch reparametrization.
\end{lemma}

\begin{proof}
The chain rule gives
\(\partial/\partial\xi^0=2\partial/\partial\eta^0\).
Thus the component vector changes from \((a,0,0,0)\) to \((2a,0,0,0)\).
For \(a=3/8\), the first presentation passes the threshold and the second
fails it.
\end{proof}

No target metric is required for this counterexample.  The repair must either
derive a source-intrinsic topology/norm, or declare a restricted presentation
group and prove the robustness statement only modulo that group.

\section{Finding F3: commonity is inserted in the principal part}

\begin{proposition}[Shared-leading-factor preload]
For any collection of sector equations
\[
 E_s(u_s)=\tau(u_s)+Z_s(u_s),
\]
where \(Z_s\) is zero order and the same nowhere-zero \(\tau\) is inserted
for every \(s\), the radical principal ideals agree:
\[
 \rad\cI_s=\langle k(\tau)\rangle.
\]
This commonity follows from the shared leading primitive alone and is
independent of every zero-order carrier map.
\end{proposition}

\begin{proof}
The principal symbol discards \(Z_s\) and equals
\(k(\tau)I\).  Its determinant is a power of \(k(\tau)\), so its radical
characteristic ideal is generated by \(k(\tau)\).  The calculation contains
no occurrence of \(M_B\).
\end{proof}

Thus the cyclic P7 carrier contributes no derivation of characteristic
commonity.  Equation \eqref{eq:ideal} is algebraically correct, but the candidate obtains
the goal property by assigning one universal leading direction to all three
sectors.  Defining \(\cC_\tau\) by equality to that same predeclared ideal
then tests consistency with the inserted primitive, not derivation of the
primitive.

The exact classification is
\[
 \texttt{new\_ontology\_primitives\_with\_repairable\_goal\_property\_preload}.
\]
This is source-side smuggling, not direct target-atlas or metric smuggling.

\section{Finding F4: the variation class excludes the decisive stress}

The attempt varies a single common field \(\tau'\) diagonally across all
sectors.  Such a variation preserves commonity by construction.

\begin{proposition}[Sector-split countermodel]
For any \(0<|\varepsilon|<1/2\), define
\[
 \tau_R=\partial_0+\varepsilon\partial_1,\qquad
 \tau_S=\partial_0,\qquad
 \tau_D=\partial_0-\varepsilon\partial_1.
\]
These fields are smooth, nowhere zero, and arbitrarily close to
\(\partial_0\) in the declared componentwise topology, but their radical
principal ideals are pairwise distinct:
\[
 \langle k_0+\varepsilon k_1\rangle,\qquad
 \langle k_0\rangle,\qquad
 \langle k_0-\varepsilon k_1\rangle.
\]
\end{proposition}

\begin{proof}
Smoothness, nonvanishing, and closeness are immediate.  Two nonzero linear
forms generate the same principal ideal over \(\R[k_0,k_1,k_2,k_3]\) only
when they differ by a nonzero scalar unit.  Comparing their \(k_0\) and
\(k_1\) coefficients shows no two displayed forms are associates when
\(\varepsilon\neq0\).
\end{proof}

This countermodel does not show that a robust common cone is impossible.  It
shows only that the current robustness claim quantifies over a class already
constrained to preserve the desired answer.  A repair must justify the
diagonal variation law from source structure or include and survive a
sector-split stress class.

\section{Finding F5: the operational composition is not typed}

The proposed sampling map has type
\[
 S_{s,0}:\Gamma(U,\R^3)\longrightarrow(\R^3)^{G_0}.
\]
The protected P7--T03 protocol response consumes a preparation,
intervention/history, guarded token, and readout context; it does not consume
an arbitrary real grid array.

\begin{proposition}[Operational interface obstruction]
The phrase “sample, then apply the P7--T03 response” is not a composition of
declared maps.  A map
\[
 B_s:(\R^3)^{G_0}\longrightarrow \mathsf{Dom}(\mathsf{P7Response}_s)
\]
or a more faithful typed span is absent.  Moreover, arbitrary real samples
need not satisfy normalization, positivity, token, history, or admissibility
conditions of the protected protocol.
\end{proposition}

\begin{proof}
Function composition requires the first codomain to lie in the second domain.
The two recorded types do not match, and no coercion or bridge is declared.
Properties required in the protocol domain are not consequences of membership
in \((\R^3)^{G_0}\).
\end{proof}

The repair must state and justify a typed bridge \(B_s\), including its source
provenance and admissibility image.  A notation-level juxtaposition is not an
operational derivation.

\section{Scoped passes}

The five failures do not erase the following valid, bounded calculations.
\begin{enumerate}[leftmargin=2.2em]
  \item The carrier matrix \(M_B\) is a three-cycle.  Among permutation
        matrices, only the identity commutes with it while fixing the named
        equipped addresses pointwise; the fixed-equipment automorphism audit
        is therefore non-erasing.
  \item Constant-copy lift followed by sampling reproduces the finite source
        data on the selected grid.
  \item Identity quotient groupoids do not collapse distinct named equipment
        inside the fixed three-control subfamily.
  \item The global generator \(p_\tau=k(\tau)\) yields overlap units
        \(u_{ij}=1\), inverses \(u_{ji}=1\), and triple-overlap cocycles
        \(u_{ij}u_{jk}u_{ki}=1\).
  \item Once the universal \(\tau\) premise is granted, the principal-symbol
        and radical-ideal calculations in \eqref{eq:ideal} are correct.
  \item The D7 interface is correctly kept separate.  No rank, adequacy, or
        activation verdict is evaluated here.
\end{enumerate}

These are internal-equipment passes.  They do not supply the missing
source-selection, all-sector, robustness, or operational-interface premises.

\section{Ten-component disposition}

\begin{longtable}{>{\raggedright\arraybackslash}p{0.06\linewidth}
                  >{\raggedright\arraybackslash}p{0.21\linewidth}
                  >{\raggedright\arraybackslash}p{0.14\linewidth}
                  >{\raggedright\arraybackslash}p{0.40\linewidth}}
\toprule
ID & Component & Verdict & Audit basis \\ \midrule
\endhead
D0 & sector family & repair & Three controls are not proved exhaustive of P7 sectors. \\
D1 & source carrier/equipment & conditional & Exact for the fixed equipped three-chain only. \\
D2 & lift and topology & repair & Lift algebra passes; the componentwise neighborhood is presentation-dependent. \\
D3a & quotient groupoids & pass & Identity quotients are non-erasing within fixed equipment. \\
D3b & overlap/source units & pass & Unit-one inverse and cocycle equations hold algebraically. \\
D4 & sector operators & repair & One universal leading field is inserted; the carrier is zero order. \\
D5a & compatibility set & repair & Equality is defined against the same predeclared shared ideal. \\
D5b & principal calculation & pass & Correct conditional on the universal-leading premise. \\
D6 & operational receipts & repair & Sampling output is not typed into the P7 response domain. \\
D7 & adequacy interface & pass/not evaluated & Separation is preserved; D7 adequacy remains outside authority. \\
\bottomrule
\end{longtable}

At the weak formula-factorization level, ten of ten slots have source-side
expressions.  Under the anchored audit, four components pass, one is
conditional, and five require repair.  Component count is not activation
credit.

\section{Verdict and repair contract}

\paragraph{Exact verdict.}
\[
 \boxed{\texttt{repair\_required\_no\_instance\_credit}}
\]
The candidate remains proposal-only repair material.  There are zero direct
target-atlas/metric imports and zero direct protected-authority imports in its
formulas, but five populated-instance defects block descriptor-instance
credit.  The D7 obstruction remains unevaluated and does not cure these
earlier defects.

\paragraph{Bounded next construction.}
After a successful checkpoint, one fresh Candidate Constructor packet may
attempt
\[
 \texttt{repair\_v22\_p4\_t02\_b2\_populated\_descriptor\_source\_intrinsic\_interfaces\_v1}.
\]
It must:
\begin{enumerate}[label=(R\arabic*),leftmargin=2.5em]
  \item type the full P7 sector family and either equip every sector or state a
        strictly restricted control result;
  \item derive a source-intrinsic topology/norm or delimit the lawful
        presentation group;
  \item derive sector leading coefficients without assuming their commonity,
        or leave universal transport explicitly unresolved;
  \item justify a variation class that includes a sector-split stress; and
  \item construct typed operational bridges \(B_s\) with admissible images.
\end{enumerate}
The packet may instead return a precise scoped obstruction.  It may not
evaluate D7, award adequacy, activate B2, or unlock P4--T03.

\paragraph{Freeze criterion.}
No family freeze is justified by this audit because the countermodels create a
materially distinct repair contract.  If the same five defects recur without
new source law or interface data, the Director should consider the scoped
negative-result label
\[
 \texttt{NDCL-V22-P4T02-B2-SHARED-TAU-SELECTOR}
\]
rather than replaying the same burden.  Such a freeze would concern the shared
transport selector strategy, not all possible B2 descriptors.

\section{Distance-to-GR accounting}

The active milestone remains \texttt{effective\_metric\_g\_eff}.  This audit
adds mathematical payload---the factorization distinction, presentation
counterexample, sector-split countermodel, shared-leading-factor proposition,
and operational type obstruction---but no positive Distance-to-GR credit.
The result sharpens the burden:
\begin{quote}
Audit every populated source primitive and map before any D7 adequacy or B2
activation decision.
\end{quote}
The exact-GR benchmark, target atlas, Lorentzian metric, causal cone,
orientation, scale, coupling, and Einstein equations remain outside this
packet.

\section{Claim ledger}

\begin{center}
\begin{tabular}{>{\raggedright\arraybackslash}p{0.48\linewidth}
                >{\raggedright\arraybackslash}p{0.38\linewidth}}
\toprule
Claim & Status \\ \midrule
All ten populated formulas weakly source-factor & draft/control pass \\
The populated descriptor is anchored in the protected source & blocked \\
Direct target geometry is imported & false for audited formulas \\
The desired commonity is preloaded by a shared source primitive & audited finding \\
The three controls cover every P7 matter sector & not established \\
The declared robustness survives source reparametrization & false for the stated norm \\
The declared robustness survives sector-split perturbations & false for the stated variation class \\
The operational composition is typed & false for the recorded maps \\
D7 adequacy is evaluated & false / not authorized \\
B2 is active or P4--T03 is unlocked & false \\
Candidate family is globally impossible & not claimed \\
\bottomrule
\end{tabular}
\end{center}

\paragraph{Terminal boundary.}
This task terminates with a repair-required audit and a guarded successor
route.  It neither adopts nor rejects ontology, makes no benchmark claim, and
takes no protected, publication, push, or external action.

\end{document}
